\documentclass[11pt,a4paper]{article}
\usepackage[margin=24mm]{geometry}
\usepackage{amsmath,amssymb,amsthm,hyperref}
\newtheorem{theorem}{Theorem}
\newtheorem{proposition}[theorem]{Proposition}
\newcommand{\F}{\mathbb F}
\newcommand{\ord}{\operatorname{ord}}
\title{Diophantine tuples from cyclic square patterns}
\author{Human Compression project: F011 application}
\date{}
\begin{document}\maketitle
This paper recovers the principal bounds of Kim--Yip--Yoo (2024),
with explicit corrections to intermediate statements and arguments.
It uses the complete companion \emph{Square patterns and cyclotomic degree
bounds over finite fields}. A Diophantine tuple consists of distinct
nonzero field elements whose products in distinct positions become squares
after adding one. Zero is allowed as the resulting square, not as a tuple
element. All fields have odd characteristic $p$. All logarithms without a
base are natural.

For a square field order $q=t^2$, its multiplicative subfield group
$\F_t^*$ consists of squares in $\F_q$: if $a\in\F_t^*$ then
$a^{(q-1)/2}=(a^{t-1})^{(t+1)/2}=1$. Products plus one of subfield
elements remain in the subfield, and zero is allowed. Thus the full
nonzero subfield is a tuple of size $\sqrt q-1$, as recorded in the
source introduction. Its cited bound $M(q)\le\sqrt q+5/2$ for all odd
$q$ is a retained background result.

\section{The principal lower bound}
\begin{proposition}[Consecutive powers]\label{powers}
Let $m\ge2$ and let $y$ be a square of order at least $m$. Define the
integer interval
\[
 J_m=\{-\lfloor m/2\rfloor,\ldots,\lceil m/2\rceil-1\}.
\]
If $q\equiv1\pmod4$, $r^2=-1$, and $y^j-1$ is a square for every
$1\le j<m$, then $\{ry^i:i\in J_m\}$ is a Diophantine $m$-tuple.
If $2$ is a square and every $y^j+1$, $1\le j<m$, is a square, then
$\{y^i:i\in J_m\}$ is a Diophantine $m$-tuple.
\end{proposition}
\begin{proof}
The interval has $m$ elements with pairwise differences at most $m-1$,
so their powers are distinct. Sums of two distinct indices have absolute
value at most $m-1$. In the first case the required expression is
$1-y^{i+j}$. For positive sums it is the product of two squares,
$-1$ and $y^{i+j}-1$. At zero it is zero. For a negative sum $-j$,
it is $(y^j-1)/y^j$, again a square. In the second case the positive
sums are covered by assumption, the zero sum gives $2$, and a negative
sum gives $(y^j+1)/y^j$. All tuple entries are nonzero.
\end{proof}

\begin{theorem}[Source Theorem 1.1, with a repaired proof]\label{main}
Let $q>7$ be a power of the odd prime $p$, with $q\equiv1,5,7\pmod8$.
Then
\[
 M(q)\ge \left\lfloor\frac p{p-1}
 \left(\frac{\frac12\log q-2\log\log q}{\log2}+1\right)\right\rfloor.
\]
For fixed $p$, along these congruence classes this is
$M(q)\ge(p/(p-1)-o(1))\log_4q$.
\end{theorem}
\begin{proof}
Let $m$ be the displayed integer. If $m\le1$ the singleton tuple suffices.
Otherwise the companion parameter lemma and square-pattern proposition
give $y$ of order at least $m$. For $q\equiv1\pmod4$ use its minus
pattern and a square root of $-1$. For $q\equiv7\pmod8$, use its plus
pattern. Here $2$ is a square: the extension degree is odd, its prime
has residue seven modulo eight, and the usual quadratic-character
supplement gives this assertion in the prime field, hence in the extension.
Proposition~\ref{powers} applies. Dividing the expression by $\log_4q$
gives the asymptotic form because $\log\log q/\log q\to0$.
\end{proof}

The exact same proof gives source Corollary~3.2, without requiring a
newly acquired proof of Cohen's theorem. It retains the classical Weil
bound, precisely as stated in the companion. This is a reworking of the
source method; it claims no historical originality. The corollary's only
additional odd prime power $q\equiv1\pmod4$ with $q\le7$ is $q=5$;
there its displayed bound is at most one, so a singleton suffices.

\begin{proposition}[The remaining residue class]
For $q>7$ with $q\equiv3\pmod8$, let $m$ be the integer in
Theorem~\ref{main}. If $m\ge2$, there is a Diophantine tuple of size
$\lfloor m/2\rfloor+1$. Consequently
\[
 M(q)\ge\left(\frac{p}{2(p-1)}-o(1)\right)\log_4q
\]
for fixed $p$ along this residue class.
\end{proposition}
\begin{proof}
The companion plus-pattern proposition does not require $2$ to be a
square. Choose the resulting $y$ and use the exponents
$0,1,\ldots,h-1$, where $h=\lfloor m/2\rfloor+1\le m$.
Two distinct exponents have sum between one and $2h-3\le m-1$.
Thus no product tests the nonsquare value $2$. Distinctness follows from
the order bound. This supplies the full asymptotic assertion of
source Remark~3.6 and does not assert maximality.
\end{proof}

\section{The stronger bound in residue class seven}
\begin{proposition}[Centered half-integer powers]\label{centered}
Let $y$ be a nonzero square of order at least $n\ge2$, and let
$z^2=y$. If $2$ and all $y^j+1$, $1\le j\le n-2$, are squares,
then
\[
 \{z^{-(n-1)},z^{-(n-3)},\ldots,z^{n-3},z^{n-1}\}
 \tag{1}
\]
is a Diophantine $n$-tuple.
\end{proposition}
\begin{proof}
Write its entries as $z^{2i-(n-1)}$, $0\le i<n$.
The ratio of two entries is $y^{i-j}$, which cannot be one for
$0<|i-j|<n$. The product of distinct entries is
$y^{i+j-(n-1)}$. Its exponent has absolute value at most $n-2$,
since the smallest distinct pair is $\{0,1\}$ and the largest is
$\{n-2,n-1\}$. A zero exponent gives $2$; positive exponents are
covered by hypothesis; and for a negative exponent $-h$ the expression
is $(y^h+1)/y^h$, a square. All entries are nonzero.
For odd $n$, (1) is the usual centered integer progression in $y$.
For even $n$, it uses half-integer powers through the chosen square root.
This is why the largest necessary exponent is still $n-2$.
\end{proof}

\begin{theorem}[Source Theorem 3.5, with a repaired construction]\label{plus-two}
If $q>7$ and $q\equiv7\pmod8$, then
\[
 M(q)\ge\left\lfloor\frac p{p-1}
 \left(\frac{\frac12\log q-2\log\log q}{\log2}+2\right)\right\rfloor.
\]
\end{theorem}
\begin{proof}
Let $n$ be the displayed integer. For $n\ge4$, the companion's
stronger parameter lemma and shorter-pattern proposition give a square
$y$ of order at least $n$ with all required $n-2$ power conditions.
The quadratic-character supplement gives that $2$ is a square, so
Proposition~\ref{centered} applies.

For completeness, every eligible field has a tuple of size three,
which handles all $n\le3$. Indeed
\[
 \#\{y\ne0,1:\chi(y)=\chi(y+1)=1\}=(q-7)/4>0.
 \tag{2}
\]
To check the identity, the complete sum of
$(1+\chi(y))(1+\chi(y+1))/4$ is $(q-1)/4$:
the linear sums vanish and $\sum\chi(y(y+1))=-1$.
The latter follows by completing the square and counting
$(u-v)(u+v)=c$ for nonzero $c$, which has $q-1$ solutions.
At zero the weight is $1/2$; at one it is one since $2$ is square;
and at $-1$ it is zero since $-1$ is nonsquare. Elsewhere the weight
is the required indicator. This proves (2).
Choose such a $y$. Squares have odd order because $(q-1)/2$ is odd,
and $y\ne1$, so its order is at least three.
The set $\{y^{-1},1,y\}$ is therefore a tuple, either directly or by
Proposition~\ref{centered}.
\end{proof}

\section{Maximal tuples}
\begin{proposition}[Source Proposition 2.3]
Every maximal Diophantine $m$-tuple satisfies $q<4^{m-1}m^2$.
\end{proposition}
\begin{proof}
For $m\ge3$, let its entries be $a_i$ and apply the companion prescribed-
character lemma with $d=2$, shifts $a_i^{-1}$ and signs $\chi(a_i)$.
Every resulting point makes all $a_i x+1$ nonzero squares. Removing the
at most $m+1$ points in $\{0,a_1,\ldots,a_m\}$ leaves at least
\[
 \frac q{2^m}-\left(\frac m2-1+2^{-m}\right)\sqrt q
 -\frac{3m}2-1.
\]
If $\sqrt q\ge2^{m-1}m$, this exceeds zero: the difference between it
and $q/2^m-m\sqrt q/2$ is at least
$(2^{m-1}-\tfrac12)m-3m/2-1=(2^{m-1}-2)m-1>0$.
It would extend the tuple, a contradiction.

For $m=2$ and $q\ge17$, use directly
$W(x)=(1+\chi(ax+1))(1+\chi(bx+1))/4$ for distinct nonzero $a,b$.
The two linear sums vanish and the quadratic sum is $-\chi(ab)$;
hence $\sum W=(q-\chi(ab))/4\ge4$. The latter identity follows, for
example, by completing the square: for nonzero $c$, counting solutions
of $v^2=u^2-c$ through $(u-v)(u+v)=c$ gives
$\sum_u\chi(u^2-c)=-1$. Removing $0,a,b$ costs at most three, so an
extending point remains. Since odd $q\ge16$ means $q\ge17$, this proves
the required strict bound when $m=2$.

A singleton is not maximal for $q\ge5$: the weight
$(1+\chi(ax+1))/2$ has total $q/2$, and removing $0,a$ costs at most
two. For $q=3$ the two nonzero elements form a tuple. The empty tuple
is never maximal. This covers every possible $m$.
\end{proof}
The asymptotic bound for all maximal tuples follows: for any fixed
$\epsilon>0$, $m\le(1-\epsilon)\log_4q$ would make
$4^{m-1}m^2<q$ for sufficiently large $q$.

\section{Exact source corrections and retained scope}
The displayed set in source Proposition~3.1 has $m+1$ terms but assumes
only $\ord(y)\ge m$. For example, over $\F_{17}$ take $m=2$, $y=-1$
and $r=4$. Both $y-1=15$ and $y^2-1=0$ are squares, and $y$ is a
square, but the displayed list is $(13,4,13)$ and has only two distinct
elements. The set is still Diophantine. Its claimed cardinality is false.
With the strict order hypothesis $\ord(y)>m$ the original proof gives
the stated $m+1$ cardinality. Proposition~\ref{powers} instead uses
exactly $m$ consecutive powers under the original weak order bound,
which is enough for the unchanged main lower bound.

Source Proposition~3.3 is not subject to that collision in its stated
section $q\equiv7\pmod8$: squares have odd order, so an even $m$ and
$\ord(y)\ge m$ imply $\ord(y)>m$. Its cyclotomic proof is valid when
the displayed quotient is interpreted as the polynomial product identity.

The introductory least-nonresidue example includes zero despite the
paper's nonzero-element convention. A valid corresponding set is
$\{1,\ldots,\lfloor\sqrt{n(p)}\rfloor\}$: products at distinct
positions plus one are less than $n(p)$, hence squares. The literal bound
$M(p)\ge\sqrt{n(p)}+1$ is false for $p=3$, where $n(3)=2$ but
$M(3)=2$. This correction leaves its stated asymptotic comparison unchanged.

The source counting proof has additional issues which the companion avoids:
the degree-one factor $\Phi_0=X$ must be counted, and the removal of zero
must be paid. The companion deliberately uses a weighted filter in general;
in the source's $q\equiv7\pmod8$ plus pattern its nonzero-point filter is
an actual indicator, since a square has odd order and cannot be a root of
any $\Phi_{2i}$ with $p\nmid i$. Source Proposition
2.3 similarly omits one point in its displayed exclusion estimate; the
proof above includes it and handles the small tuple sizes separately.

Source Theorem~3.5, in the section $q\equiv7\pmod8$, replaces the final
$+1$ in Theorem~\ref{main} by $+2$. This is a separate, stronger claim.
Its proof asserts $T\le B+2$ for
$m=\lfloor\alpha(B+2)\rfloor$, which does not follow and can fail.
For $q=7^{11}$ the prescribed integer is $m=10$, giving $T=9$, whereas
$8.6006<B+2<8.6008$. There is also an exact small-parameter contradiction
to the displayed character-count lower bound. For $q\equiv7\pmod8$
and $m=2$, its count is
\[
 N(2)=\#\{y\ne0,1:\chi(y)=\chi(y+1)=1\}=(q-7)/4.
\]
Indeed the complete weight sum is $(q-1)/4$ by the elementary quadratic
sum evaluated in the preceding proof; zero has weight $1/2$, one has
weight one, and $-1$ has weight zero. The source instead gives
$N(2)\ge(q+\sqrt q)/4-1$, since $T=2$ and $S=1$, which is larger
than the exact count. The value $q=7^5$ lies in the theorem's branch
and gives its prescribed $m=2$, so this is an error inside the proof's
actual parameter range. This refutes that intermediate estimate, not
the claimed existence of a tuple.
Its final logarithmic deduction also does not follow merely from
$\alpha<2$. These defects in the printed route do not refute the
existence theorem. Theorem~\ref{plus-two} proves the exact stronger bound
using centered half-integer powers, the shorter cyclotomic filter, an
explicitly stronger order exclusion, and the corrected parameter bound
$T\le B+2+1/p$. All these steps are written in full in the two papers.

The source conjecture $M(q)\ge(1-o(1))\log_2q$ is recorded without an
answer or a current-status claim. The cited square-field upper bound,
integer-tuple history and GRH least-nonresidue estimates remain attributed
background, not fresh theorems proved here. No new Lean or numerical
compression measurement is supplied.
\end{document}
